\subsection{単調で有界な数列}
数列が\textbf{上に有界}であるとは,すべての $n$ について $a_n\le M$ を満たす一定の実数 $M$ が存在することをいう.下に有界であることも,同様に一定の下限によって定義する.増加を続けていても,上に有界なら無限に大きくなることはできない.

\begin{theorem}{単調有界数列の収束：採用する事実}
単調増加する数列が上に有界なら,有限の極限値へ収束する.単調減少する数列が下に有界なら,同じく有限の極限値へ収束する.
\end{theorem}

この定理の完全な証明には,実数の連続性に関する性質が必要になる.本書では実数の基礎理論までさかのぼらず,この定理を採用する.ただし,適用するときに単調性と有界性の両方を確かめることは省略しない.収束すると予想した式に極限値を代入するだけでは,収束の証明にはならないからである.



